\textbf{Proof.}
Put \(d_0=n-1\), \(p=\dim(S)\), \(\beta=p/d_0\), and \(s=\rho^2\). Every \(A\in\mathcal C_n\) is a convex combination of matrices \(zz^\top\) with \(z\in\mathcal Z_n\). Consequently,
\[
A\succeq0,
\qquad
A\mathbf 1_n=0,
\qquad
\operatorname{diag}(A)=(1,\ldots,1)^\top,
\qquad
\operatorname{tr}(A\mid H_n)=n.
\tag{1}
\]
Since \(H_n\) has dimension \(d_0\), the trace bound gives
\[
\lambda_{\max}(A\mid H_n)\geq \frac{n}{d_0}=c_n.
\tag{2}
\]
Complete randomization has covariance \(c_nP_H\), so its normalized loss equals \(c_n\) for every quality level.

At perfect quality, the feasible unit vectors are exactly the unit vectors in \(S\). Thus
\[
q_1(A;S)=\lambda_{\max}(P_SAP_S\mid S),
\qquad
v^*(1;S)=\kappa(S).
\tag{3}
\]
If strict improvement occurs at some \(s<1\), compactness supplies a minimizing covariance \(A\), and monotonicity of the feasible spherical shells gives
\[
\kappa(S)\leq q_1(A;S)\leq q_{\sqrt{s}}(A;S)=v^*(\sqrt{s};S)<c_n.
\tag{4}
\]
Conversely, suppose \(\kappa(S)<c_n\), and let \(A_0\in\mathcal C_n\) attain \(\kappa(S)\). If \(s_j\uparrow1\) and \(u_j\) maximizes the spherical-shell loss at \(s_j\), compactness of the unit sphere in \(H_n\) gives a convergent subsequence. Its limit has score energy one, hence lies in \(S\). It follows that
\[
q_{\sqrt{s}}(A_0;S)\longrightarrow q_1(A_0;S)=\kappa(S)
\qquad(s\uparrow1).
\tag{5}
\]
Therefore strict improvement occurs at some \(s<1\). Equations \((3)\)--\((5)\) prove
\[
S\in\mathfrak A_n
\quad\Longleftrightarrow\quad
\kappa(S)<c_n.
\tag{6}
\]

By \(\text{thm:cone-lift-dual}\), and because the value at multiplier \(\eta=0\) is at least \(c_n\) by \((2)\),
\[
\begin{aligned}
v^*(\sqrt{s};S)<c_n
&\Longleftrightarrow
\text{there exist }A\in\mathcal C_n\text{ and }\eta>0\text{ such that}\\
&\hspace{2.2cm}
\lambda_{\max}(A+\eta P_S\mid H_n)-\eta s<c_n\\
&\Longleftrightarrow
s>\inf_{A\in\mathcal C_n,\,\eta>0}
\frac{\lambda_{\max}(A+\eta P_S\mid H_n)-c_n}{\eta}\\
&\Longleftrightarrow s>\gamma(S).
\end{aligned}
\tag{7}
\]
The middle equivalence uses the defining property of an infimum: a strict upper bound on the infimum is exceeded by some member of the indexed set. Taking \(A=c_nP_H\) shows \(\gamma(S)\leq1\), because
\[
\lambda_{\max}(c_nP_H+\eta P_S\mid H_n)=c_n+\eta.
\tag{8}
\]
Thus the set of improving squared qualities is \((\gamma(S),1]\) when \(\gamma(S)<1\), and is empty when \(\gamma(S)=1\). Equations \((6)\)--\((8)\) prove the three-way equivalence and
\[
\rho^*(S)=\sqrt{\gamma(S)}
\quad\text{on }\mathfrak A_n,
\qquad
\rho^*(S)=1
\quad\text{off }\mathfrak A_n.
\tag{9}
\]

We next prove the dimension bound and its exact equality case. Taking the trace on \(H_n\) gives, for every feasible \(A\) and every \(\eta>0\),
\[
\lambda_{\max}(A+\eta P_S\mid H_n)
\geq \frac{n+\eta p}{d_0}
=c_n+\eta\beta.
\tag{10}
\]
Therefore \(\gamma(S)\geq\beta\).

For sufficiency of equal leverage, we establish rather than assume the required geometry of the balanced covariance hull. Let
\[
\mathcal L
=\{B\in\mathbb S^n:\operatorname{diag}(B)=0,\ B\mathbf1_n=0\}.
\tag{11}
\]
We claim that
\[
\operatorname{aff}(\mathcal C_n)=c_nP_H+\mathcal L.
\tag{12}
\]
The inclusion from left to right follows from \((1)\). For the reverse inclusion, put
\[
\mathcal V=\operatorname{span}\{zz^\top-c_nP_H:z\in\mathcal Z_n\}\subseteq\mathcal L.
\]
Suppose a symmetric matrix \(Q=(q_{ij})\) is orthogonal to \(\mathcal V\). Then \(z^\top Qz\) is constant over all balanced sign vectors. For a set \(T\subseteq\{1,\ldots,n\}\) of size \(n/2\), let \(z_T=2\mathbf1_T-\mathbf1_n\). Given four distinct indices \(i,j,k,\ell\), choose a set \(U\) of size \(n/2-2\) disjoint from them; this is possible for every even \(n\geq4\). The alternating sum of the four equal quantities indexed by
\[
U\cup\{i,k\},\quad U\cup\{j,k\},\quad
U\cup\{i,\ell\},\quad U\cup\{j,\ell\}
\]
is zero. Expanding \(z_T^\top Qz_T\) cancels all constant and linear terms and yields
\[
q_{ik}-q_{jk}-q_{i\ell}+q_{j\ell}=0
\qquad(i,j,k,\ell\text{ distinct}).
\tag{13}
\]
These four-point identities imply that there are scalars \(a_1,\ldots,a_n\) such that
\[
q_{ij}=a_i+a_j\qquad(i\neq j).
\tag{14}
\]
For completeness, choose three distinct anchors \(r,s,t\), set
\[
a_r=\frac{q_{rs}+q_{rt}-q_{st}}2,
\quad
a_s=\frac{q_{rs}+q_{st}-q_{rt}}2,
\quad
a_t=\frac{q_{rt}+q_{st}-q_{rs}}2,
\]
and, for every other \(i\), set \(a_i=q_{ir}-a_r\). Applying \((13)\) to \((i,t,r,s)\) and \((i,s,r,t)\) gives \(q_{is}=a_i+a_s\) and \(q_{it}=a_i+a_t\). Applying it to \((i,r,j,s)\) then gives \(q_{ij}=a_i+a_j\) for two non-anchor indices. This proves \((14)\), including the case \(n=4\).

By \((14)\), the difference between \(Q\) and \(\mathbf1_na^\top+a\mathbf1_n^\top\) is diagonal. Both a diagonal matrix and \(\mathbf1_na^\top+a\mathbf1_n^\top\) are orthogonal to \(\mathcal L\). Hence \(\mathcal V^\perp\subseteq\mathcal L^\perp\). The reverse inclusion follows from \(\mathcal V\subseteq\mathcal L\), so \(\mathcal V=\mathcal L\), proving \((12)\).

We also prove the relative-interior fact used below. Enumerate the distinct matrices \(z z^\top\) as \(K_1,\ldots,K_m\). Their uniform average is \(c_nP_H\): permutation symmetry makes all off-diagonal entries equal, and annihilation of \(\mathbf1_n\) makes that common entry \(-1/(n-1)\). If \(B\in\mathcal L\), equation \((12)\) permits an expression
\[
B=\sum_{j=1}^m b_j(K_j-c_nP_H)
=\sum_{j=1}^m d_jK_j,
\qquad \sum_{j=1}^m d_j=0.
\]
For all sufficiently small \(|t|\), every coefficient \(m^{-1}+td_j\) is nonnegative, and hence
\[
c_nP_H+tB=\sum_{j=1}^m(m^{-1}+td_j)K_j\in\mathcal C_n.
\tag{15}
\]
Thus \(c_nP_H\) lies in the relative interior of \(\mathcal C_n\) inside the affine space in \((12)\).

Now assume equal leverage, so \((P_S)_{ii}=p/n\) for every \(i\). Since \((P_H)_{ii}=d_0/n\), the matrix
\[
D=\beta P_H-P_S
\]
has zero diagonal; it also annihilates \(\mathbf1_n\), because \(S\subseteq H_n\). Therefore \(D\in\mathcal L\). By \((15)\), for all sufficiently small \(t>0\),
\[
A_t=c_nP_H+tD\in\mathcal C_n.
\]
For the multiplier \(\eta=t\),
\[
A_t+tP_S=(c_n+t\beta)P_H,
\]
so its quotient in the definition of \(\gamma(S)\) equals \(\beta\). Together with \((10)\), this proves \(\gamma(S)=\beta\).

For necessity, suppose \(\gamma(S)=\beta\). Choose a minimizing sequence \((A_k,\eta_k)\) with \(A_k\in\mathcal C_n\), \(\eta_k>0\), and
\[
g_k:=\frac{\lambda_{\max}(A_k+\eta_kP_S\mid H_n)-c_n}{\eta_k}
\longrightarrow\beta.
\tag{16}
\]
Pass to a subsequence such that \(\eta_k\) converges in \([0,\infty]\). Put
\[
B_k=A_k+\eta_kP_S,
\qquad
\overline\lambda_k=c_n+\eta_k\beta,
\qquad
\Delta_k=\lambda_{\max}(B_k\mid H_n)-\overline\lambda_k
=\eta_k(g_k-\beta)\geq0.
\tag{17}
\]

If \(\eta_k\to0\), let \(E_k=B_k-\overline\lambda_kP_H\) on \(H_n\). Its eigenvalues sum to zero and each is at most \(\Delta_k\); hence each is at least \(-(d_0-1)\Delta_k\). Therefore
\[
\|E_k\|_{\mathrm{op}}\leq(d_0-1)\Delta_k=o(\eta_k).
\tag{18}
\]
Using the unit diagonal of \(A_k\), the \(i\)-th diagonal entry of
\[
A_k=c_nP_H+\eta_k(\beta P_H-P_S)+E_k
\]
gives, after subtracting \(c_n(P_H)_{ii}=1\), dividing by \(\eta_k\), and applying \((18)\),
\[
0=\beta\frac{d_0}{n}-(P_S)_{ii}+o(1)
=\frac pn-(P_S)_{ii}+o(1).
\]
Thus \((P_S)_{ii}=p/n\) for every \(i\).

If \(\eta_k\to\eta_0\in(0,\infty)\), compactness of \(\mathcal C_n\) permits a further subsequence with \(A_k\to A_0\). Equations \((16)\)--\((17)\) give \(\Delta_k\to0\). The limiting symmetric operator has largest eigenvalue equal to its average eigenvalue, so every one of its \(d_0\) eigenvalues equals that average:
\[
A_0+\eta_0P_S=(c_n+\eta_0\beta)P_H.
\tag{19}
\]
Taking diagonal entries in \((19)\), and again using \(\operatorname{diag}(A_0)=1\), yields \((P_S)_{ii}=p/n\) for every \(i\).

Finally, if \(\eta_k\to\infty\), positive semidefiniteness and trace \(n\) imply \(\|A_k\|_{\mathrm{op}}\leq n\). Hence
\[
\lambda_{\max}\!\left(P_S+\eta_k^{-1}A_k\mid H_n\right)
=\eta_k^{-1}c_n+g_k
\longrightarrow\beta.
\]
The left side converges to \(\lambda_{\max}(P_S\mid H_n)=1\), because \(p\geq1\). Thus \(\beta=1\), so \(p=d_0\), \(S=H_n\), and \(P_S=P_H\) has diagonal \(d_0/n=p/n\). These three cases exhaust every minimizing sequence and prove
\[
\gamma(S)=\frac{p}{n-1}
\quad\Longleftrightarrow\quad
\operatorname{diag}(P_S)=\frac pn(1,\ldots,1)^\top.
\tag{20}
\]

The consequences for \(\rho^*(S)\) follow from \((9)\), \((10)\), and \(p\leq n-1\). If \(S\notin\mathfrak A_n\), then \(\gamma(S)=1\) by the established equivalence and \(0\leq\gamma(S)\leq1\), and the convention \(\rho^*(S)=1\) preserves the same dimension lower bound. The strict-improvement implication follows directly from \((7)\), and scaling by \(4M/n>0\) converts the normalized comparison into the displayed comparison of \(V^*\) and \(V_{\mathrm{CRD}}\).

It remains to justify computability. Enumerate one representative \(z^{(1)},\ldots,z^{(m)}\) from each balanced sign-pair and write
\[
A(a)=\sum_{j=1}^m a_jz^{(j)}z^{(j)\top},
\qquad a\in\Delta_m.
\]
If \(Q_S\) is an orthonormal basis matrix for \(S\), then \(P_S=Q_SQ_S^\top\), and
\[
\kappa(S)=\min_{a\in\Delta_m,\,t}t
\quad\text{subject to}\quad
tI_p-Q_S^\top A(a)Q_S\succeq0.
\tag{21}
\]
For fixed \(s\), equation \((7)\) is the semidefinite feasibility test
\[
a\in\Delta_m,quad \eta>0,qquad
(c_n+\eta s)I_{d_0}
-Q_H^\top\{A(a)+\eta P_S\}Q_H\succ0,
\tag{22}
\]
where \(Q_H\) is any orthonormal basis matrix for \(H_n\). The test is monotone in \(s\); one-dimensional bisection, with semidefinite margin optimization at each query and non-strict bounds at the limiting boundary, computes \(\gamma(S)\) and \(\rho^*(S)\) to any prescribed tolerance. The finite saddle program in \(\text{thm:exact-oracle-saddle}\) and its simplex dual compute an \(\varepsilon\)-optimal covariance and sign-pair weights, whose sign symmetrization is an attaining design in \(\mathcal D_n\). Finally, \((20)\) is checked directly by forming \(Q_SQ_S^\top\) and inspecting its diagonal. Every input to \((21)\)--\((22)\) is determined by \(n\) and the pre-randomization score basis. \(\square\)